\subsection{Path second moments and the mixed central estimate}\label{proof:R2-C044}
\paragraph{R44.1 The local central estimates}

All expectations use the finite measure and kernel of
Section~\ref{sec:prelim}. This module proves the path second moment
and the mixed estimate used at a center whose residual begins with
a unit arm.

At each branching vertex, maximally pair equal side-arm lengths. A
residual list consists of the remaining distinct lengths in increasing
order. Each paired arm is still directed and retained. The central
families treated here are exactly
\fitdisplay{
(1,s),\quad s\ge2,
}{R44.1.2}
and
\fitdisplay{
(1,s,a,b,u_1,v_1,\ldots,u_k,v_k),\qquad
 2\le s<a<b<u_1<v_1<\cdots<u_k<v_k,\quad k\ge0.
}{R44.1.3}
Together they are all nonempty even residual lists of minimum one.

For every \(s\ge2\), the actual outward path message satisfies
\[
  E f_s^2\le1.
\]
For the central residual in (R44.1.3), set
\[
  Q=W_a\bar f_b\prod_{i=1}^k W_{u_i}\bar f_{v_i}.
\]
The same-label mixed estimate is
\[
  E[f_s^2Q^8]\le e^{-17d/64}.
\]
For (R44.1.2), take \(Q=1\) and use the path second moment.
The strict ordering in (R44.1.3) includes \(s=2\) and \(a=3\).
These are the local inputs to Table~\ref{tab:five-center}; the
compound reference-second construction is proved in
Proposition~\ref{prop:five-compound}.

\paragraph{R44.2 Original finite calculus and genuine path moments}

All operators are on the same original real L2(mu). Delete only
zero-mass host atoms, which leaves every original integral unchanged. Put
\fitdisplay{
B=-K,\quad u=B\one,\quad T=J+B,\quad A=J-B=T^*,\quad
 p=1-u,\quad q=1+u,\quad Jv=(Ev)1,
}{}
\fitdisplay{
f_n=A^{n}\one,\quad \bar f_n=T^{n}\one,\quad
 h_n=Tf_{n-1},\quad W_n=A\bar f_{n-1},\quad f_0=\bar f_0=1.
}{R44.2.1}
An \(f_n\) arm points entirely away from its root; \(\bar f_n\) points entirely
toward it; \(h_n\) has its first edge toward the root and every remaining
edge away along the arm; \(W_n\) has its first edge away from the root and
every remaining edge toward that first neighbor. In particular
\(f_1=W_1=p\) and \(\bar f_1=h_1=q.\) An incoming connector of length L has
operator \(T^L,\allowbreak \) with every original edge directed toward the receiver.

Write
\fitdisplay{
\delta_j=\|B^{j}\one\|_2^2,\quad d=\delta_1,\quad
 \theta=13/32,\quad \kappa=19/32,\quad \tau=d-\delta_2.
}{}
The exact C014 finite calculus gives
\fitdisplay{
0\le d\le1/3,\quad d^2\le \delta_2\le\theta d,\quad
 \delta_{j+1}\le\theta \delta_j,\quad \|B\|^2\le\theta,\quad
 |BV(x)|\le\|V\|_2,\quad Eu=0,
}{R44.2.2}
\fitdisplay{
EB^{2j}1=(-1)^j\delta_j,\qquad EB^{2j+1}1=0.
}{R44.2.3}
In particular \(\delta_2\) is a spectral moment, whereas \(d^2\) is the square of
the first moment; they are never identified.

For \(G_n(z)=\sum_{j=0}^{n-1}z^j\), with \(G_0=0,\allowbreak \) define
\fitdisplay{
\beta_1=\beta_2=0,\qquad
 \beta_n=\sum_{j=1}^{\lfloor(n-1)/2\rfloor}(-1)^{j-1}\delta_j
       \quad(n\ge3).
}{}
The exact finite positive convolution and its complete reverse are
\fitdisplay{
S_n^+=1+G_n(B)u
       =\bar f_n+\sum_{j=3}^n\beta_j\bar f_{n-j},
}{}
\fitdisplay{
S_n^-=1-G_n(-B)u
       =f_n+\sum_{j=3}^n\beta_jf_{n-j}.
}{R44.2.4}
Every coefficient \(\beta_j,\allowbreak  j\ge 3,\allowbreak \) lies in [tau, d], and \(\beta_3=\beta_4=d.\)
The sums are finite and empty at \(n=1,\allowbreak \) 2. Their means are
\fitdisplay{
ES_n^+=ES_n^-=1-\beta_{n+1}.
}{R44.2.5}
The positive path messages \(f_n,\allowbreak  \bar f_n\) lie in [0, 2] and have means at most
one. The mean assertion is not extended to \(h_n\) or \(W_n: Eh_2=EW_2=1+d.\)

The finite geometric estimate needed below is
\fitdisplay{
\|G_n(B)\|\le1\quad(n=0,1\text{ or }n\ge3),\qquad
 \|G_2(B)\|^2\le1+\theta.
}{R44.2.6}
Here is the finite calculation, including every integer n. On an
original real skew rotation block put \(z=i \sqrt{x}\), \(0\le x\le \theta<1/2.\)
The exact finite geometric identity gives
\fitdisplay{
|G_n(z)|^2=\frac{1-2\operatorname{Re}(z^n)+x^n}{1+x}.
}{}
For odd \(n\ge 3,\allowbreak  Re(z^n)=0\) and \(x^n\le x.\) For n divisible by four,
\(x^n-2 x^{n/2}\)\ensuremath{\le}\(0\le x.\) For n congruent to two modulo four and
\(n\ge 6,\allowbreak \)
\fitdisplay{
x^n+2x^{n/2}\le x^6+2x^3
       =x(x^5+2x^2)\le(17/32)x\le x.
}{}
The estimate uses only \(x\le 1/2\); it therefore covers every original
rotation block and the zero block. The remaining exact cases are
\(G_0=0,\allowbreak  G_1=1\) and \(|G_2(z)|^2=1+x.\) Orthogonal summation over
the original finite blocks proves (R44.2.6), with no spectral limit.

Write \(G_n(z)=E_n(x)+zD_n(x).\) For every \(n\ge 1,\allowbreak \)
\fitdisplay{
0\le E_n\le1,\quad 0\le D_n\le1;\qquad E_1=1,\ D_1=0.
}{R44.2.7}
For \(n\ge 2\) the two coordinates lie in [1-x, 1]. These are finite
alternating partial sums. No restriction \(x\ge d\) is made.

Let nu be the original spectral measure of -\(B^2\) at u. Its mass is d,
its support is [0, theta], and int \(x^j d\nu= \delta_{j+1}\). For a real
polynomial \(Q(z)=E(x)+zD(x),\allowbreak \)
\fitdisplay{
\|Q(B)u\|^2=\int(E^2+xD^2)d\nu,\qquad
 E[Q(B)u]=-\int D\,d\nu,\qquad
 EB[Q(B)u]=-\int E\,d\nu.
}{R44.2.8}
Every constant-vector cross below is computed on this same original
space before any norm estimate.

\paragraph{R44.2.1 Complete ordinary second moments of the actual reference arms}

For \(s=2\) the exact actual formula is
\fitdisplay{
f_2=1-u+Bu,\quad Ef_2=1-d,\quad
 Ef_2^2=1-d+\delta_2\le1-\kappa d.
}{R44.2.9}
Indeed, \(\E Bu=-d\) and \(\langle u,Bu\rangle=0\).

For every \(s\ge 3,\allowbreak \) (R44.2.4) gives \(0\le f_s\le S_s^-\) and (R44.2.6) gives
\fitdisplay{
E(S_s^--1)^2=\|G_s(-B)u\|^2\le d.
}{}
Consequently
\fitdisplay{
Ef_s^2\le E(S_s^-)^2
  =1-2\beta_{s+1}+\|G_s(-B)u\|^2
  \le1-d+2\delta_2\le1-3d/16.
}{R44.2.10}
Thus \(Ef_s^2\le 1\) for every genuine \(s\ge 2,\allowbreak \) directly from the
finite G and positive-convolution calculus. Only these finite bounds
enter this estimate.
The unit arm is different: \(Ep^2=1+d.\)

\paragraph{R44.2.2 A uniform fourth bound for every genuine \(s\ge 2\)}\label{part:R44.2.2}

For -\(1\le y\le 1,\allowbreak \)
\fitdisplay{
(1+y)^4\le1+4y+11y^2,
}{R44.2.11}
because the difference is \(y^2(1-y)(5+y)\ge 0.\)

At \(s=2,\allowbreak  f_2\) in [0, 2], \(E(f_2-1)=-d\) and
\(E(f_2-1)^2=d+\delta_2.\) Therefore
\fitdisplay{
Ef_2^4\le1+7d+11\delta_2
       \le1+\frac{367}{32}d
       \le e^{367d/32}.
}{R44.2.12}

At \(s\ge 3,\allowbreak \) put \(g_s=min(2,\allowbreak  S_s^-).\) This is only an analytical
nonnegative majorant: the actual arm remains \(f_s.\) Positivity in (R44.2.4)
and \(f_s\le 2\) give \(0\le f_s\le g_s\le 2.\) Moreover
\fitdisplay{
E(g_s-1)\le-\beta_{s+1}\le0,\qquad
 E(g_s-1)^2\le E(S_s^--1)^2\le d.
}{}
The second inequality is pointwise before integration: clipping \(S_s^-\)
at two can only decrease its squared distance from one. Using (R44.2.11),
\fitdisplay{
Ef_s^4\le Eg_s^4\le1-4\beta_{s+1}+11d
       \le1+11d\le e^{367d/32}.
}{R44.2.13}
This proves the uniform actual-path estimate
\fitdisplay{
\boxed{Ef_s^4\le e^{367d/32}\quad(s\ge2).}
}{R44.2.14}
It holds for every finite s, including the genuine \(s=3/G_3\) case.
There is no large-length limit and no arbitrary positive-function
replacement of \(f_s.\)

\paragraph{R44.2.3 Equal original pairs}

For each actual equal-length pair,
\fitdisplay{
f_\ell+h_\ell=2Ef_{\ell-1},\qquad
 0\le f_\ell h_\ell\le(Ef_{\ell-1})^2\le1.
}{R44.2.15}
Complete reversal gives \(0\le \bar f_\ell W_\ell\le 1.\) Products of these actual
factors are denoted \(E_v\) at their original vertices. They can be compared
with one only inside a nonnegative original integral or positive
propagation. 

\paragraph{R44.3 The full signed exponential estimate at exponent sixteen}

On the original probability space, apply
\hyperref[part:Q81.2]{the signed-exponential theorem} with
\(W=K(-V)=BV\), \(U=\|V\|_2\), and \(s=t\).
The bounds in (R44.2.2) give \(|W|\le U\) and
\(EW^2=\|BV\|_2^2\le\theta U^2\).
Thus every real complete vector \(V\), with either sign of its actual
mean \(E(BV)\), satisfies
\fitdisplay{
E e^{tBV}\le
 \exp\{tEBV+(17/64)t^2\|V\|^2\},\qquad
 V\text{ \text{any } real original vector},\quad t\ge0.
}{R44.3.1}
Taking \(t=16\) gives \((17/64)16^2=68\), hence
\fitdisplay{
\boxed{Ee^{16BV}\le e^{16EBV+68\|V\|^2}.}
}{R44.3.4}

\paragraph{R44.4 New central mixed moment, with every short and finite-pair case}\label{part:R44.4}

Let the genuine central lengths satisfy (R44.1.3), and put
\fitdisplay{
Q=W_a\bar f_b\prod_{i=1}^kW_{u_i}\bar f_{v_i}.
}{R44.4.1}
This directs each of the original central arms other than 1, s
exactly once. The central unit arm will not be used anywhere in this
section.

Central lemma. For the same actual \(f_s\) and Q in (R44.4.1),
\fitdisplay{
\boxed{E[f_s^2Q^8]\le e^{-17d/64}.}
}{R44.4.2}
Arbitrary actual equal-pair factors may be restored by positivity,
though the final assembly assigns all central equal pairs to its
separate C106 component.

\paragraph{R44.4.1 The finite mask and the higher-pair envelope}

For \(4\le c<e\) the full reversal of \eqref{eq:L90:pair} gives
\fitdisplay{
W_c\bar f_{e}\le\exp\{BV_{c,e}-16d/5\},\qquad
 V_{c,e}=\sum_{j=c+1}^{e}B^{j-2}u.
}{R44.4.3}
This is an actual original pair bound, also explicitly proved in
C096 \mbox{\hyperref[part:R44.2.2]{part R44.2.2}} and preserved under complete reversal. Only this
exact pair interface is needed from C186's larger terminal theorem.

For any finite polynomial P with coefficients zero or one, the four
residue classes of its powers at i \(\sqrt{x}\) give
\fitdisplay{
|P(i\sqrt x)|^2\le\frac{1+x}{(1-x^2)^2}
 \le\frac{32768}{16245}<\frac{101}{50}=:C_*.
}{R44.4.4}
For each coordinate, bound its positive and negative subcollections
and take the larger absolute bound; do not add the two. Consequently
a finite union of disjoint exponent intervals beginning at degree h
has original vector norm squared at most \(C_*\theta^h\) d.
All intervals used below are disjoint because the original arm
lengths are strictly increasing.

If V is such a union beginning at degree at least three, its positive
terms in EBV begin at \(\delta_4.\) Hence
\fitdisplay{
EBV\le\frac{\delta_4}{1-\theta^2}<\frac d{12},\qquad
 \|V\|^2\le C_*\theta^3d<\frac9{64}d.
}{R44.4.5}
The first rational comparison is
\(\theta^3/(1-\theta^2)=2197/27360<1/12.\)
The norm comparison is \(221897<230400\) after the common denominator
1638400. These bounds do not grow with the number of intervals.

\paragraph{R44.4.2 Every \(a\ge 4,\allowbreak \) with every finite later pair block}

There are \(r=k+1\) actual pairs. Let V be their single full finite mask.
Equations (R44.4.3)-(R44.4.5), followed by (R44.3.4), give
\fitdisplay{
EQ^{16}\le
 \exp\left[-\left\{\frac{256r}{5}-\frac43-\frac{153}{16}\right\}d\right].
}{R44.4.6}
At \(r=1\) the coefficient is \(9673/240>40.\) Each further original
pair adds \(256/5\) to this coefficient, while the full-vector estimate
already pays every cross term. Thus, by Cauchy on the original mu
and the actual path bound (R44.2.14),
\fitdisplay{
E[f_s^2Q^8]\le(Ef_s^4)^{1/2}(EQ^{16})^{1/2}
 \le e^{-(40-367/32)d/2}
 \le e^{-17d/64}.
}{R44.4.7}


\paragraph{R44.4.3 The actual \(a=3\) envelope, derived before any moment estimate}

Here the strict order forces \(s=2.\) Write \(n=b-3\ge 1\) and
\fitdisplay{
V_0=B^2G_n(B)u,\qquad \Gamma=3d-2\delta_2=d+2\tau.
}{}
The exact \(W_3+\bar f_3=2E\bar f_2=2-2d\) and (R44.2.4) give
\fitdisplay{
W_3+\bar f_b=2-d+\sum_{j=4}^{b}B^{j}\one
                 -\sum_{j=3}^{b}\beta_j\bar f_{b-j}.
}{}
Retain the two genuine last terms \(\beta_b \bar f_0\) and
\(\beta_{b-1}\bar f_1,\allowbreak \) both with coefficient at least tau. Since \(b\ge 4,\allowbreak \)
both indices are at least three. Their sum is at least tau(2+u).
Therefore
\fitdisplay{
W_3+\bar f_b\le2+BV_0-\tau u-\Gamma.
}{}
The left side is the sum of two actual nonnegative messages.
AM-GM and \(1+t\le e^t\) imply
\fitdisplay{
Q_0:=W_3\bar f_b\le
       \exp\{B(V_0-\tau1)-\Gamma\}.
}{R44.4.8}
The vector is \(V_0\) minus tau times the constant vector 1. It is
not \(V_0-\tau\) u.

In (R44.2.7), write \(G_n=E_n+zD_n.\) Then \(V_0\) has coordinates
-\(xE_n-zxD_n.\) From (R44.2.8),
\fitdisplay{
0\le EV_0=\int xD_n\,d\nu\le \delta_2,\qquad
 EBV_0=\int xE_n\,d\nu\le \delta_2.
}{R44.4.9}
Thus the full constant square is
\fitdisplay{
\|V_0-\tau1\|^2=\|V_0\|^2+\tau^2-2\tau EV_0.
}{R44.4.10}
Every following estimate starts with (R44.4.10). Its cross is not
discarded when needed in the genuine \(n=2\) exception.

\paragraph{R44.4.4 No later pairs, low domain \(0<d\le 1/16:\) n not equal to two}

If \(n=1\) or \(n\ge 3,\allowbreak \) (R44.2.6) gives \(\lVert V_0\rVert ^2\le \delta_3.\)
By (R44.4.9)-(R44.4.10),
\fitdisplay{
\|V_0-\tau1\|^2\le\theta \delta_2+\tau^2.
}{}
The signed mean in (R44.3.4) satisfies
\(EB(V_0-tau1)=EBV_0\le \delta_2,\allowbreak \) since \(Eu=0.\) Put \(t=\delta_2/d\) only
after these original integrated inequalities. Then
\fitdisplay{
\frac{\log EQ_0^{16}}d
 \le-48+(48+68\theta)t+68d(1-t)^2.
}{R44.4.11}
For every fixed \(d\le 1/16\) its derivative in t is at least
\(48+68\theta-136/16>0\) on [0, theta]. Its maximum is therefore
at \(t=\theta.\) The remaining d coefficient is nonnegative, so use
\(d=1/16.\) The resulting loss is exactly
\fitdisplay{
48\kappa-68\theta^2-\frac{17}{4}\kappa^2
       =\frac{64631}{4096}>15.
}{R44.4.12}
This handles $b=4$ and every $b\ge6$.

\paragraph{R44.4.5 No later pairs, low domain: the genuine \(n=2,\allowbreak  b=5\) exception}

Now \(G_2=1+B,\allowbreak \) so it must not be treated as a contraction:
\fitdisplay{
V_0=B^2u+B^3u,\quad EV_0=\delta_2,\quad
 \|V_0\|^2=\delta_3+\delta_4.
}{}
The cross between consecutive B powers vanishes by skewness.
The nonzero constant cross in (R44.4.10) is exactly -2tau \(\delta_2.\)
Hence
\fitdisplay{
\|V_0-\tau1\|^2
 \le\theta(1+\theta)\delta_2+(d-\delta_2)(d-3\delta_2).
}{R44.4.13}
This expression comes from the same actual vector; a potentially
negative second summand is kept with its sign. Equations (R44.3.4)
and (R44.4.8) give
\fitdisplay{
\frac{\log EQ_0^{16}}d
 \le-48+\{48+68\theta(1+\theta)\}t
                   +68d(1-t)(1-3t).
}{R44.4.14}
For fixed \(d\le 1/16\) the derivative in t is at least
\(48+68\theta(1+\theta)-17>0.\) At \(t=\theta\) the last term is
nonpositive, because \(1-3\theta=-7/32.\) Dropping only that
nonpositive term yields the loss
\fitdisplay{
48\kappa-68\theta^2(1+\theta)
       =\frac{104187}{8192}>12.
}{R44.4.15}
Consequently (R44.4.11)-(R44.4.15) prove, for every \(b\ge 4,\allowbreak \)
\fitdisplay{
EQ_0^{16}\le e^{-12d},\qquad 0<d\le1/16.
}{R44.4.16}
Since the reference is the actual \(f_2,\allowbreak \) Cauchy and (R44.2.12) give
\fitdisplay{
E[f_2^2Q_0^8]\le
 e^{-(12-367/32)d/2}=e^{-17d/64}.
}{R44.4.17}

\paragraph{R44.4.6 No later pairs, high domain \(1/16\le d\le 1/3\)}

The actual direction (R44.4.8) is unchanged. From (R44.2.6), (R44.4.9)
and (R44.4.10),
\fitdisplay{
\|V_0-\tau1\|^2
 \le(1+\theta)\delta_3+\tau^2
 \le(1+\theta)\theta^2d+d^2.
}{R44.4.18}
Also \(\Gamma\ge 35d/16.\) For \(d\ge 1/16,\allowbreak \)
\fitdisplay{
(1+\theta)\theta^2=\frac{7605}{32768}
 <\frac{7752}{32768}
 =\frac{969}{256}\frac1{16}
 \le\frac{969}{256}d.
}{}
Adding \(d^2\) in (R44.4.18) therefore gives
\fitdisplay{
\|V_0-\tau1\|^2\le(35d/16)^2\le\Gamma^2.
}{}
The original row bound in (R44.2.2), applied to the full real vector,
shows \(B(V_0-tau1)\le \Gamma\) pointwise. Thus
\fitdisplay{
Q_0\le1,\qquad
 E[f_2^2Q_0^8]\le Ef_2^2\le1-\kappa d
                         \le e^{-\kappa d}
                         \le e^{-17d/64}.
}{R44.4.19}
This pointwise conclusion is proved only on its stated high
domain.  Both proofs include \(d=1/16\) and use the same
literal \(W_3 \bar f_b\) direction.

\paragraph{R44.4.7 \(a=3\) with every nonempty finite later pair block}

Let \(k\ge 1.\) The original order gives \(u_1>b\ge 4,\allowbreak \) hence \(u_1\ge 5.\)
Let
\fitdisplay{
V_*=\sum_{i=1}^k\sum_{j=u_i+1}^{v_i}B^{j-2}u,\qquad
 V=V_0+V_*.
}{}
The entire set of powers in V is one disjoint finite zero-one
mask beginning at degree two. Therefore
\fitdisplay{
\|V\|^2\le C_*\theta^2d
             =\frac{17069}{51200}d<\frac9{25}d.
}{R44.4.20}
For the later vector alone, the positive terms in \(EBV_*\) begin
at \(\delta_4,\allowbreak \) while the negative terms in \(EV_*\) begin at -\(\delta_3.\)
Thus, for the actual finite subcollections,
\fitdisplay{
EBV_*\le\frac{\delta_4}{1-\theta^2}<d/12,\qquad
 EV_*\ge-\frac{\delta_3}{1-\theta^2}>-d/5.
}{R44.4.21}
The second rational comparison is \(\theta^2/(1-\theta^2)\)
=\(169/855<1/5.\) Combined with \(EV_0\ge 0,\allowbreak \) these imply
\fitdisplay{
\|V-\tau1\|^2
 =\|V\|^2+\tau^2-2\tau EV
 \le\frac9{25}d+\frac75d^2
 \le\frac{62}{75}d.
}{R44.4.22}
The bound includes the root-pair cross terms and the constant-vector cross term.

Using (R44.4.3), (R44.4.8) and (R44.4.21),
\fitdisplay{
Q\le\exp\{B(V-\tau1)-\Gamma-16kd/5\},
}{}
\fitdisplay{
\Gamma+16kd/5-EB(V-\tau1)
 \ge\{57/32+16k/5-1/12\}d.
}{R44.4.23}
At \(k=1,\allowbreak \) sixteen times the last coefficient is \(2351/30.\)
The entire variance cost in (R44.3.4) is at most \(68(62/75)=4216/75.\)
Their exact difference is
\fitdisplay{
\frac{2351}{30}-\frac{4216}{75}
       =\frac{3323}{150}>22.
}{}
Every further original pair adds \(256/5.\) Hence, on the whole
original domain without a d split,
\fitdisplay{
EQ^{16}\le
 \exp[-\{3323/150+256(k-1)/5\}d]\le e^{-22d}.
}{R44.4.24}
Cauchy with (R44.2.12) again proves (R44.4.2), with substantial room.
There is no upper bound on k or on any original length.

Parts R44.4.2 and 4.3-4.7 treat \(a\ge 4\) and \(a=3,\allowbreak \) respectively.
At \(d=0\) the actual messages all equal one, by leaf integration. This completes the new central lemma.


